\documentclass[xcolor=table]{beamer}
\usepackage{beamerthemesplit}
\usepackage{wrapfig}
\usetheme{SPbGU}
\usepackage{pdfpages}
\usepackage{amsmath}
\usepackage{cmap} 
\usepackage[T2A]{fontenc} 
\usepackage[utf8]{inputenc}
\usepackage[english]{babel}
\usepackage{indentfirst}
\usepackage{amsmath}
\usepackage{tikz}
\usepackage{multirow}
\usepackage[noend]{algpseudocode}
\usepackage{algorithm}
\usepackage{algorithmicx}
\usepackage{fancyvrb}
\usetikzlibrary{shapes,arrows}
%usepackage{fancyvrb}
%\usepackage{minted}
%\usepackage{verbments}

\newtheorem{mytheorem}{Theorem}
\renewcommand{\thealgorithm}{}


\beamertemplatenavigationsymbolsempty

\title[F\# OpenCL C TP]{F\# OpenCL C Type Provider}
%\subtitle[YaccConstructor]{Parsing techniques for graph analysis}
% То, что в квадратных скобках, отображается в левом нижнем углу. 
\institute[JetBrains Research]{
JetBrains Research, Programming Languages and Tools Lab  \\
Saint Petersburg University
}

% То, что в квадратных скобках, отображается в левом нижнем углу.
\author[Semyon Grigorev]{Kirill Smirenko, \textbf{Semyon Grigorev}}

\date{September 27, 2018}

\definecolor{orange}{RGB}{179,36,31}

\begin{document}
{
\begin{frame}[fragile]
  \begin{tabular}{p{3.5cm} p{5.5cm} p{1cm}}
   \begin{center}
      \includegraphics[height=1.5cm]{pictures/jetbrainsResearch.pdf}
    \end{center}
    &
    \begin{center}
      \includegraphics[height=1.5cm]{pictures/icfplogo.png}
    \end{center}
    &
    \begin{center}
      \includegraphics[height=1.5cm]{pictures/SPbGU_Logo.png}
    \end{center} 
  \end{tabular}
  \titlepage
\end{frame}
}

\begin{frame}[fragile]
  \transwipe[direction=90]
  \frametitle{GPGPU}
  \begin{minipage}[m]{0.40\linewidth}
\raisebox{-0.5\totalheight}{\includegraphics[height=0.85\textheight]{pictures/cpu_gpu.png}}
\end{minipage}\hfill
\begin{minipage}[m]{0.6\linewidth}
General purpose computations on grphical processor units
\begin{itemize}
      \item (Almost) SIMD architecture 
      \item Huge amount of ``simple'' ALUs on single chip
      \item May be a good choice for huge data processing
\end{itemize}

\end{minipage}

\end{frame}

\begin{frame}[fragile]
  \transwipe[direction=90]
  \frametitle{General purpose applications of GPGPU}
  \begin{itemize}
    \item Initially is sientific computations
    \begin{itemize}
      \item Phis
      \item Math
      \item Chem
    \end{itemize}
    \item But more amd more general application
      \begin{itemize} 
        \item Finance/Banking
        \item Bioinformatics
        \item Data Analytics and Data Science (Hadoop, Spark ...)
        \item Security analytics (log processing)
      \end{itemize}
  \end{itemize}
\end{frame}

\begin{frame}[fragile]
  \transwipe[direction=90]
  \frametitle{Problem: GPGPU <-> High level programming}
  \begin{itemize}
   \item .NET, JVM, etc
   \item Interaction is a problem!
  \end{itemize}
\end{frame}

\begin{frame}[fragile]
  \transwipe[direction=90]
  \frametitle{Existing solutions and problems}
  \begin{itemize}
  \item Generative approach (haskell, Alea, etc): good byt what about code reusing
  \item Simple drivers (textual, for example)
  \end{itemize}
\end{frame}

\begin{frame}
  \transwipe[direction=90]
  \frametitle{Type providers}
  our focus is .NET, so it is the way to solve our problem
\end{frame}

\begin{frame}
  \transwipe[direction=90]
  \frametitle{Type providers}
  more inside, exaples of R, COM, INI providers
  %https://fsprojects.github.io/FSharp.Configuration/index.html
  \begin{itemize}
    \item Config \\
     {\includegraphics[width=0.85\textwidth]{pictures/YamlConfigProvider.png}}
  \end{itemize}
\end{frame}

\begin{frame}[fragile]
  \transwipe[direction=90]
  \frametitle{Brahma.FSharp}
  \begin{itemize}
  \item F\# quotations to OpenCL C translator
  \item Runtime
      \begin{itemize}
        \item Comand queue
        \item Context management
        \item Memory management
        \item F\# aliases for OpenCL-specifc functions
      \end{itemize}

  \end{itemize}
\end{frame}

\begin{frame}
  \transwipe[direction=90]
  \frametitle{OpenCL C type provider}
\begin{itemize}
\item Imorove OpenCL C lexer, parser and  translator
\item Unify kernels on client side
\item Improve user exp
\end{itemize}

\end{frame}
            
\begin{frame}
  \transwipe[direction=90]
  \frametitle{Architecture}
  Diagram
\end{frame}

\begin{frame}
  \transwipe[direction=90]
  \frametitle{Limitations }
\begin{itemize}
\item Only (small) subset of OpenCL C
 \begin{itemize}
 \item h files
 \item preprocessor
 \item subset of sintax
 \item !!!
 \end{itemize}
\item Very simple type mapping
\item !!!
\item !!!
\end{itemize}

\end{frame}

            
\begin{frame}
  \transwipe[direction=90]
  \frametitle{Examples}         
  Screens
\end{frame}     

\begin{frame}
  \transwipe[direction=90]
  \frametitle{Future work}         
\begin{itemize}
\item Imorove OpenCL C lexer, parser and  translator
\item Unify kernels on client side
\item Improve user exp
\end{itemize}
\end{frame}     
            
            
\begin{frame}
  \transwipe[direction=90]
  \frametitle{Summary}         
\begin{itemize}
\item F\# OpenCL C type provider
\item Prototype
\item Type-safe using of existing OpenCL C code in F\# applications
\end{itemize}
\end{frame}           
            
\begin{frame}
\transwipe[direction=90]
\frametitle{Contact Information}
\begin{itemize}
  \item Semyon Grigorev: \href{mailto:s.v.grigoriev@spbu.ru}{s.v.grigoriev@spbu.ru}
  \item Kirill Smirenko: \href{mailto:k.smirenko@gmail.com}{k.smirenko@gmail.com}
\end{itemize}
\begin{itemize}
  \item Brahma.FSharp: \href{https://github.com/YaccConstructor/Brahma.FSharp}{https://github.com/YaccConstructor/Brahma.FSharp}
\end{itemize}
\hspace{2cm}
\center{\huge{Thakns!}}
\end{frame}
\end{document}
